% Template report — compile test for sustech-labreport.cls
% Compile with xelatex (Times New Roman + xeCJK for the Chinese name).
% Prefer scripts/auto_gen_template.py over copying this file by hand: the
% script fills the cover from the live CLI (course, name, SID, experiment
% date). Replace every <...> marker that remains.
% Section names follow the course convention — keep them, they auto-number.
\documentclass[12pt]{sustech-labreport}
% The course rulebook writes figure captions as "Fig. 1"; the class alone
% prints "Figure 1". Keep this line. Table captions stay "Table N".
\captionsetup[figure]{name=Fig.}
\begin{document}

\makesutitle
  {Experiments for Advanced Materials Science and Engineering I}
  {Experiment N: <Experiment title>}
  {<Name>}
  {<SID>}
  {YYYY/MM/DD}
  {}
  {}

% Headings are auto-numbered (1., 2., ...) by the class — never type the
% numbers by hand and never use starred headings. Subheadings below a
% section auto-number 5.1, 5.1.1.

\section{Objectives}
\begin{enumerate}
  \item To <goal 1 of the experiment, taken from the manual>.
  \item To <goal 2: the characterisation or skill to master>.
\end{enumerate}

\section{Principles}
<Condensed theory behind the experiment — key formulas and the principle
figure. Keep within one page. Do not teach textbook background.>

% No blank line before display math (it creates a spurious paragraph break),
% and the label goes on the SAME line as \begin{equation}:
\begin{equation}\label{eq:main}
<the governing relation>
\end{equation}

\section{Experimental Materials and Setups}
Chemicals: <reagent 1>, <reagent 2>, ...
Apparatuses: <apparatus 1>, <apparatus 2>, ...
Tools: <tool 1>, <tool 2>, ...

% APPARATUS / SETUP PHOTOS BELONG HERE (or in Procedures) — never in Results.

\section{Experimental Procedures}
% In-class parameters override the manual's; write what was actually done.
\begin{enumerate}
  \item <Step 1 of the in-class operation>.
  \item <Step 2>.
  \item <Step 3>.
\end{enumerate}

\section{Results and Discussion}
% THREE sub-sections, in the order of the work: you see, then you calculate,
% then you analyze. Phenomena comes first and is plain prose, no numbers.
\subsection{Phenomena}
<What was seen: color, appearance, when it changed, what was surprising.>

\subsection{Experimental Data}
<2-3 sentences of prose carrying the key measured numbers, then the table.>

% TABLE RECIPE — three-line booktabs; caption ABOVE; label right after the
% caption. Every cell names the quantity and its unit, in words:
%   "Mass of \ce{CuCl2.2H2O} (g)"   NOT   "m(CuCl2.2H2O) / g"
% NO Note block by default — add one only for a condition that cannot be
% expressed in the table and changes how a value must be read.
\begin{center}
  \captionof{table}{<Concise title of the table.>}\label{tab:props}
  \begin{tabular}{lc}
    \toprule
    Quantity & Value \\
    \midrule
    Density of sample A ($\rho$, g\,cm$^{-3}$)  & 1.25 \\
    Tensile strength of sample A ($\sigma$, MPa) & 85.2 \\
    \bottomrule
  \end{tabular}
\end{center}
% Only a LARGE table uses a real float instead — \begin{table}[htbp] with the
% caption still ABOVE the tabular. Columns are centred ({cc}, {ccc}).
% =============================================================================

% =============================================================================
% FIGURE RECIPE
%
% 1. ONE THEME PER FIGURE. A before/after or sample-1/2/3 comparison never
%    shares a figure with an apparatus photo — split them into two figures.
% 2. Subfigures: \quad between items in a row, \\ (or \bigskip) between rows.
%    NEVER \hfill — it stretches the row to the margins and leaves visible
%    gaps whenever the image widths do not sum to \linewidth.
% 3. Sub-captions are SHORT — two or three words. A sentence-long
%    sub-caption wraps into a tall ragged column next to the picture.
% 4. Each \subfloat needs BOTH the short [...] sub-caption AND a \label
%    inside its {} group.
% 5. The main caption carries the description and cites the letters with
%    \protect\subref{...} — never type (a)/(b) by hand.
% 6. Caption BELOW the figure, \label on the next line.
% 7. Every figure is discussed in prose ("as shown in Fig.~\ref{...}").
%
% Single picture:
%   \begin{figure}[htbp]
%     \centering
%     \includegraphics[height=6cm,keepaspectratio]{fig/result.jpg}
%     \caption{<What the reader sees and the take-away.>}
%     \label{fig:result}
%   \end{figure}
%
% Two pictures that answer one question together:
%   \begin{figure}[htbp]
%     \centering
%     \subfloat[Before heating]{\includegraphics[height=4cm,keepaspectratio]{fig/before.jpg}\label{fig:before}}\quad
%     \subfloat[After heating]{\includegraphics[height=4cm,keepaspectratio]{fig/after.jpg}\label{fig:after}}
%     \caption{Pictures of the sample before \protect\subref{fig:before} and
%     after \protect\subref{fig:after} heating.}
%     \label{fig:heating}
%   \end{figure}
%
% Four pictures, two rows (\quad inside a row, \\ between rows):
%   \subfloat[Room temperature]{\includegraphics[height=4cm]{fig/a}\label{fig:a}}\quad
%   \subfloat[Above transition]{\includegraphics[height=4cm]{fig/b}\label{fig:b}}\\[1ex]
%   \subfloat[Apparatus]{\includegraphics[height=4cm]{fig/c}\label{fig:c}}\quad
%   \subfloat[Product]{\includegraphics[height=4cm]{fig/d}\label{fig:d}}
%
% Multi-panel captions follow the published style — letters inline, one
% clause per letter:
%   \caption{Cross-sectional images ... for 5 min at 250 $^\circ$C:
%   (a) no additive, (b) 1 ppm, (c) 2 ppm, (d) 5 ppm.}
%
% SIZING: height= for photos sharing a row (equal heights keep it tidy),
% width= for a lone image. Never set width and height without
% keepaspectratio — it stretches the picture.
% MICROGRAPHS AND PHOTOS NEED THE SCALE BAR AND THE MAGNIFICATION stated in
% the caption (the teacher requires both). If the user did not give one,
% ASK — do not infer it from the filename or from looking at the image.
% =============================================================================

\subsection{Error Analysis}
<Required by the rubric: quantitative or qualitative error sources, and how
 they were reduced. Never invent an uncertainty the manual or the user did not
 state.>

\section{Conclusion}
<1-2 short paragraphs, ~300 words: what was successfully prepared, the
 critical numbers against theory, the cause of the shortfall. No personal
 pronouns, no "we learned".>

\section{Questions and Further Thoughts}
\begin{enumerate}
  \item <Restate the manual's question.>

  <Answer in 2-3 sentences: one physical reason plus one observation from this
  experiment.>
\end{enumerate}

% REFERENCES — only if real literature is cited. NEVER cite the course manual;
% if there is no other literature, delete the section entirely.

\newpage
\section{Raw Data and Preview Report}
% The scanned handwritten preview and the raw-data record, appended as pages.
% Never OCR them and never use their content as report content.
\begin{center}
  \includegraphics[width=0.86\textwidth]{fig/preview-1.png}
\end{center}
\begin{center}
  \includegraphics[width=0.86\textwidth]{fig/rawdata-1.png}
\end{center}

\end{document}
